\chapter*{Abstract}

NextGenPB is a finite-element-based linear Poisson-Boltzmann equation solver focused on computational efficiency, which is geared toward high-performance computing, but does not yet exploit GPUs to enhance performance even more.
In this thesis, I used various techniques to accelerate it with Nvidia's CUDA API.
The focus lay on the solution of the linear system, as well as the energy calculation.
A few additional host-side optimizations decreased the overall time even further while also improving memory efficiency.

For the linear equation I used Nvidia's AMGX linear solver library with great effect, achieving speedups of up to 10x.
I created two configurations, one focused on memory efficiency and the other on speed; both massively outperformed LIS.
I found that AMGX scales well across multiple GPUs on a single node but not across multiple nodes, where the communication overhead nullified or even reversed the gains.

For the energy calculation I created a naive implementation and a fast multipole method.
The naive implementation brought a huge speedup of 24.6x on an RTX 3080 and 321x on an A100 for the energy calculation of the 1VSZ molecule.
On top of that, the fast multipole method offered even greater speedup for compute-bound problems, improving the speedup on the RTX 3080 to 117x, all while maintaining the original accuracy of NextGenPB.
Even though the fast multipole method loses its advantage for non-compute-bound problems, this only applies to problems that take very little time to begin with.

These optimizations made it possible to compute the entire H1N1 capsid on a single node with 4 H200 GPUs in under one hour.
They also make it possible to compute medium-sized problems using fewer cluster nodes.
